Training produces a set of adapters $\{\Delta_c\}_{c \in \mathcal{C}^*}$, one per identified capability. Jointly merging all adapters into a single set of weights leads to parameter interference: capabilities impose conflicting optimization pressures on shared parameters, and strengthening one degrades others---a well-documented challenge in multi-task learning~\citep{yu2020gradient, fifty2021efficiently} and continual fine-tuning~\citep{kirkpatrick2017overcoming, french1999catastrophic}. We instead compose adapters at inference time, selecting and weighting them based on the demands of each task instance.

\paragraph{Retrieval-augmented routing.}
We construct a corpus of interaction trajectories from each capability's training environment, where each trajectory demonstrates how a capability was correctly applied to solve a task instance. At inference time, we embed the initial task prompt $o_1$ with an embedding model and retrieve the top-$k$ most similar trajectories from the corpus via cosine similarity. The capabilities whose trajectories appear among these top-$k$ form the candidate set for classification. This reduces the number of options the classifier must distinguish among, which becomes important as the total number of capabilities grows.

We then use the base model $\pi_\theta$ itself as a zero-shot classifier over the candidates. Each candidate is assigned a distinct label token (\texttt{A}, \texttt{B}, \texttt{C}, \ldots) and presented with its natural-language description and the retrieved trajectory showing its correct application in a similar scenario. An additional label represents the unmodified base model. From this prompt, $\pi_\theta$ produces next-token probabilities over the label set. Because routing requires only coarse discrimination among the filtered candidates---not solving the underlying task---the base model serves as an effective zero-shot classifier without additional training (Section~\ref{sec:ablations}).

\paragraph{Adapter selection.}
The router's label probabilities define the adapter weights:
\[
\alpha_c \;\propto\; \pi_\theta(c \mid o_1), \qquad c \in \mathcal{C}^*,\quad \alpha_c \geq 0.
\]
If the base-model label receives the highest probability, all adapter weights are zero and inference proceeds with $\pi_\theta$ unmodified. Otherwise, we apply nucleus (top-$p$) filtering: capabilities are sorted by descending $\alpha_c$, and we retain the smallest set whose cumulative probability exceeds a threshold $p$, zeroing the rest and re-normalizing. This allows the number of active adapters to vary per instance: when the router is confident, a single adapter dominates; when the task spans multiple capabilities, multiple adapters are retained.

\paragraph{Adapter composition.}
Given the routing weights $\alpha_c$, we merge the selected adapters directly into the base model parameters before inference. Each adapter $\Delta_c$ consists of low-rank factors $B_c \in \mathbb{R}^{d_{\text{out}} \times r}$ and $A_c \in \mathbb{R}^{r \times d_{\text{in}}}$. For each layer, we precompute the weighted sum of adapter deltas and add it to the base weights:
\begin{equation}
W' \;=\; W \;+\; \sum_{c}\, \alpha_c \, B_c A_c,
\label{eq:compose}
\end{equation}
where $W$ denotes the original base-model weights. Since each adapter contributes an additive low-rank term, precomputing $W'$ is mathematically equivalent to applying the adapters independently during the forward pass---no approximation is introduced. Inference then runs on $W'$ through the standard forward pass with no per-token overhead from adapter composition. The first merge incurs a one-time cost to cache the original weights $W$ in CPU memory. Each subsequent swap restores $W$ from this cache via a bulk transfer and applies the new adapter deltas, requiring only one rank-$r$ multiplication per adapted layer per active adapter. The weights $\alpha_c$ are determined once from the initial task prompt and held fixed throughout the interaction, so this merge is performed once per task instance.
